\chapter{Introduction}

In 2024, Mexico received more than 67 billion current US dollars in remittances, making it the second-largest recipient country in the world after India \parencite{worldbank_remittances}. This figures represents approximately 3.7\% of Mexico's GDP (reference?), highlighting the macroeconomic significance of these transfers. With roughly 96\% of remittances to Mexico originating from the United States (reference?), and the average Mexican household in the US sending approximately US\$12,800 per year \parencite{isidro_martinez_2025}, a natural question arises: does an economic shock in the US propagate to Mexico through the remittance channel? Specifically, do Mexican migrants absorb the shock and maintain their transfer levels constant, or do they reduce remittance outflows in response to a decline in their current income or wealth? The answer has direct implications for understanding the vulnerability of Mexican households to external shocks and the extent to which the remittance channel acts as a shock transmission mechanism.

\par

The closest paper to ours is \textcite{caballero_2023}, who study how the Great Recession affected remittance receipt among Mexican households. Using a shift-share research design that exploits variation in migration networks across Mexican municipalities, they find that households more connected to severely hit U.S. labor markets experiencing larger employment declines were significantly less likely to receive remittances. This confirms that the remittance channel acts as a transmission mechanism for US labor demand shocks. However, their analysis is limited to the extensive margin (i.e. whether the households receive remittances at all), while not providing any estimate for the intensive margin (i.e. how much remittance flows decline in response to the shock). Quantifying this magnitude is central to our paper, as it determines the actual income loss borne by Mexican households. Moreover, the Great Recession was a global shock of exceptional magnitude, which simultaneusly contracted both the U.S. and Mexican economies, with Mexico recording a 6.6 points decline in real GDP in 2009 \parencite{villarreal_2010}. This makes it difficult to isolate the effect of the U.S. economic decline on Mexican remittance inflows. Furthermore, shocks of this magnitude are uncommon, which limits the external validity of the findings derived from them.  

\par

We address these challenges by exploiting the localized and more frequent nature of natural disasters. From 1980 to 2024, the U.S. has sustained 403 weather and climate disasters for which the individual damage costs reached or exceeded US\$1 billion, with cumulative costs that exceeds US\$2.915 trillion \parencite{ncei_billions_2025}. These events generate localized reductions in firm productivity, leading to lower family income \parencite{natural_disasters_US} and therefore plausibly reducing the capacity of affected migrant households to remit.

\par 

The relationship between natural disasters and remittances has been studied before, but existing work focuses on the opposite side of the corridor, analyzing what happens to bilateral remittance flows when the \textit{receiving} country is struck. \textcite{bettin_2023} study the evolution of remittance flows from 103 Italian provinces to 79 developing countries and find that remittances increase by approximately 1.8\% on impact following a disaster in the recipient country, rising to about 2.4\% four months later and rapidly decreasing after one year. This provides a useful benchmark for the speed and persistence of remittance adjustment, although it operates through a different mechanism than ours, where the shock originates on the \textit{sender} side.

\par

We contribute to this literature by providing causal evidence on how a localized US shock propagates to Mexico through the remittance channel. To this end, we exploit Hurricane Ian (September 2022), the costliest hurricane in Florida's history, which caused US\$\,109.5 billion in damage \parencite{nhc_ian_2026}, as a natural experiment. 
The storm struck the southwestern coast of Florida but did not affect Mexico, providing a clean one-sided shock to the sender side of the remittance corridor. We estimate bilateral remittance flows between each US state and each Mexican municipality by applying a migration weighting matrix to Banxico data on remittance inflows by U.S. State of origin \parencite{banxico_ce168}. Then, we implement a difference-in-differences design comparing Florida-linked municipalities against those connected to unaffected states. We find that ... Our findings suggest that ..., implying that ...